% ============================ FRONT MATTER ===========================
\chapter*{How to read this book}
\markboth{How to read this book}{}
\addcontentsline{toc}{chapter}{How to read this book}

This book follows the five AMATH301 lectures in order. Every chapter starts from the
pictures, then states the definitions and theorems exactly as they appear in the course,
then proves them in full. Before each proof you will find a short paragraph that says,
in plain words, \emph{why} the argument works. If you only remember that paragraph on exam
day, you can usually rebuild the proof from it.

\section*{The boxes you will meet}

\begin{tcolorbox}[corners=plum,colback=plumtint,bpgridbg=lavgrey!70,title=Definition]
Definitions are framed in plum on a faint grid. Learn them word for word: most exam
questions start by testing whether you know exactly what is being asked.
\end{tcolorbox}

\begin{tcolorbox}[corners=coral,colback=coraltint,bpgridbg=coral!10,title=Theorem]
Theorems are framed in coral. Each one is followed by its proof.
\end{tcolorbox}

\begin{intuition}
The blue box gives the idea behind a statement or a proof, with no formulas if possible.
\end{intuition}

\begin{trap}
Coral warning boxes point to the mistakes that cost the most marks.
\end{trap}

\begin{reallife}{}
Navy boxes connect the mathematics to physics, engineering, economics or machine learning.
\end{reallife}

The tinted margin on the right is for short reminders and small pictures, and it is
also the space where you can write your own notes when you print the book. Exercises at
the end of each chapter are rated
{\color{coral}$\bullet$}{\color{lavgrey}$\bullet\bullet$} (direct application),
{\color{coral}$\bullet\bullet$}{\color{lavgrey}$\bullet$} (exam standard) and
{\color{coral}$\bullet\bullet\bullet$} (challenge). Full solutions are in
Appendix~\ref{app:sol}, and a formula sheet closes the book.

\clearpage
\section*{The map of the course}
\markboth{The map of the course}{}

Multivariate calculus is one long story about a single idea: \emph{near a point, a smooth
function looks like a linear one}. The diagram below shows how the five lectures build on
each other. Read it from top to bottom.

\begin{center}
\begin{tikzpicture}[
  lec/.style={draw=plum, line width=0.9pt, fill=plumtint,
     text width=4.8cm, align=left, inner sep=6pt, font=\small},
  head/.style={fill=plum, text=white, font=\small\bfseries\scshape, inner sep=4pt,
     anchor=south west},
  arr/.style={-{Stealth[length=2.6mm]}, coral, line width=1.1pt},
  arrd/.style={-{Stealth[length=2.2mm]}, cyanbp, line width=1pt, dashed},
  lab/.style={font=\scriptsize, text=coraldark, fill=white, inner sep=2pt},
]
\node[lec] (L1) at (-3.1,0) {\textbf{Functions of several variables}\\[2pt]
  domain, graph, level curves\\ open sets, partial derivatives\\
  gradient $\nabla f$, differentiability};
\node[head] at (L1.north west) {Lecture I};
\node[lec] (L2) at (3.1,0) {\textbf{Chain rule and gradient}\\[2pt]
  $\frac{d}{dt}f(X(t))=\nabla f\cdot X'(t)$\\ tangent planes, $D_Af$\\
  steepest ascent, energy};
\node[head] at (L2.north west) {Lecture II};
\node[lec] (L3) at (-3.1,-3.6) {\textbf{Higher derivatives}\\[2pt]
  $D_1D_2f=D_2D_1f$\\ differential operators\\
  Taylor's formula with $(H\cdot\nabla)^r$};
\node[head] at (L3.north west) {Lecture III};
\node[lec] (L5) at (3.1,-3.6) {\textbf{Derivative as a linear map}\\[2pt]
  Jacobian matrix $J_F$, differentials\\ chain rule $(G\circ F)'=G'F'$\\
  inverse and implicit functions};
\node[head] at (L5.north west) {Lecture V};
\node[lec, text width=6cm] (L4) at (0,-7.4) {\textbf{Maxima and minima}\\[2pt]
  critical points $\nabla f(P)=O$, quadratic form\\ boundary, closed and bounded sets\\
  Lagrange multipliers $\nabla f=\lambda\nabla g$};
\node[head] at (L4.north west) {Lecture IV};

\draw[arr] (L1.east) -- node[lab, above=3pt]{compose} (L2.west);
\draw[arr] (L1.south) -- node[lab, right=2pt]{differentiate again} ([yshift=0.45cm]L3.north);
\draw[arr] (L2.south) -- node[lab, left=2pt]{vector-valued maps} ([yshift=0.45cm]L5.north);
\draw[arr] (L3.south) -- ++(0,-0.75) -| node[lab, pos=0.25, above]{quadratic term}
   ([xshift=-1.6cm,yshift=0.45cm]L4.north);
\draw[arrd] (L5.south) -- ++(0,-0.75) -| node[lab, text=cyanbp, pos=0.25, above]{constraint surfaces}
   ([xshift=1.6cm,yshift=0.45cm]L4.north);
\draw[arr] (L2.east) -- ++(0.3,0) |- node[lab, pos=0.78, above]{$\nabla f=O$ at extrema} (L4.east);
\end{tikzpicture}
\end{center}

\begin{keyidea}
Every chapter answers one question about the approximation
\[
 f(P+H)\;\approx\; f(P) \;+\; \underbrace{\nabla f(P)\cdot H}_{\text{Lectures I, II, V}}
 \;+\; \underbrace{\tfrac12\,(H\cdot\nabla)^2 f(P)}_{\text{Lectures III, IV}} .
\]
The linear term tells us slopes, tangent planes and directions of fastest change. The
quadratic term tells us whether a flat point is a hill top, a valley or a saddle.
\end{keyidea}
